% SPDX-License-Identifier: CC-BY-4.0
\documentclass[11pt,letterpaper]{article}

\usepackage[T1]{fontenc}
\usepackage{amsmath,amssymb}
\usepackage{tikz}
\usepackage{pgfplots}
\pgfplotsset{compat=1.18}
\usepackage{algorithm}
\usepackage{algpseudocode}
\usepackage{booktabs}
\usepackage{enumitem}
\usepackage[hidelinks]{hyperref}

\newcommand{\N}{\mathbb{N}}
\newcommand{\Z}{\mathbb{Z}}
\newcommand{\Q}{\mathbb{Q}}
\newcommand{\R}{\mathbb{R}}
\newcommand{\C}{\mathbb{C}}
\newcommand{\Map}[3]{#1 \colon #2 \to #3}
\newcommand{\Mod}[1]{#1\text{-}\mathrm{Mod}}
\newcommand{\GroupHom}{\mathrm{Hom}_{\mathbf{Grp}}}
\newcommand{\RingHom}{\mathrm{Hom}_{\mathbf{Ring}}}

\newtheorem{theorem}{Theorem}
\newtheorem{lemma}[theorem]{Lemma}
\newtheorem{corollary}[theorem]{Corollary}
\newtheorem{definition}[theorem]{Definition}
\newtheorem{remark}[theorem]{Remark}
\newenvironment{proof}{\noindent\textit{Proof.}}{\hfill\(\square\)}

\title{The even row of a local split}
\author{ascendnoosphere}
\date{\today}

\begin{document}
\maketitle

\section*{The even row}

Let \(T\) be the classical shortcut on the natural numbers: \(T(n)=n/2\) when \(n\) is even, and \(T(n)=(3n+1)/2\) when \(n\) is odd. A parity word is read chronologically. The letter \(0\) is one even step and the letter \(1\) is one odd step. The empty word fixes its input.

The even row consists of the integers \(n>1\) that are divisible by \(2\). Its source word is the single letter \(0\), and its parent word is empty. The parent attached to such an \(n\) is \(n/2\). This note proves that row and nothing else.

\begin{definition}
\label{def:words}
Write \(S_0(n)=n/2\) and \(S_1(n)=(3n+1)/2\). For a parity word \(W\), the evaluation \(R_W(n)\) is \(n\) when \(W\) is empty, and \(R_{W'}(S_b(n))\) when \(W=bW'\). The even source word is \(W_s=0\), and the even parent word is \(W_p=\varepsilon\).
\end{definition}

\begin{lemma}
\label{lem:identity}
If \(2\) divides \(n\), then \(2\cdot(n/2)=n\) and \(T(n)=n/2\). The source word and the parent word then satisfy
\[
R_{W_s}(n)=n/2, \qquad R_{W_p}(n/2)=n/2, \qquad 2\cdot R_{W_s}(n)=n.
\]
\end{lemma}

\begin{proof}
Division with remainder gives \(n=(n/2)\cdot 2+(n\bmod 2)\). The hypothesis \(n\bmod 2=0\) removes the remainder, so \(n=(n/2)\cdot 2\). Commutation of the factors is \(2\cdot(n/2)=n\).

The even branch of \(T\) is \(n/2\), so \(T(n)=n/2\). The word \(W_s=0\) applies \(S_0\) once and then the empty word, hence \(R_{W_s}(n)=n/2\). The empty word fixes \(n/2\), so \(R_{W_p}(n/2)=n/2\). Multiplying the source evaluation by \(2\) and using the doubling identity gives \(2\cdot R_{W_s}(n)=n\).
\end{proof}

\begin{theorem}
\label{thm:even-split}
Let \(n>1\) be even, and set \(p=n/2\). Then \(p\) is a positive integer, \(p<n\), \(T(n)=p\), and
\[
R_{W_s}(n)=R_{W_p}(p).
\]
Thus one even letter carries \(n\) to a strictly smaller positive parent, and the empty word on that parent is already the common endpoint.
\end{theorem}

\begin{proof}
Lemma~\ref{lem:identity} gives \(2p=n\), \(T(n)=p\), \(R_{W_s}(n)=p\), and \(R_{W_p}(p)=p\). The two word evaluations are therefore equal.

The quotient \(p\) cannot be \(0\): otherwise \(n=2\cdot 0=0\), contradicting \(n>1\). Since \(p\ge 1\),
\[
p<p+p=2p=n.
\]
Hence \(0<p<n\).
\end{proof}

\begin{theorem}
\label{thm:odd-outside}
If \(n\) is odd, then \(n\) is not a member of the even row. In particular \(3>1\) lies outside the row, and \(2\cdot(3/2)=2\neq 3\).
\end{theorem}

\begin{proof}
Membership requires \(n\) even. An odd residue is not the residue \(0\) modulo \(2\), so the membership conjunction fails.

The integer \(3\) has residue \(1\) modulo \(2\), so the preceding sentence applies. Separately, integer division gives \(3/2=1\), and \(2\cdot 1=2\). The positive difference \(3-2=1\) shows that the doubling identity of Lemma~\ref{lem:identity} fails at \(3\).
\end{proof}

\begin{remark}
Theorem~\ref{thm:even-split} is one local row. It does not treat an input congruent to \(2\) modulo \(3\), an input congruent to \(1\) modulo \(4\), or an input of the form \(2^a u-1\) with \(u\) odd. It does not classify a complement of this row, and it does not assert that every integer greater than \(1\) has a parent in any family. Theorem~\ref{thm:odd-outside} is only the obstruction to reading the even row as a cover.
\end{remark}

\end{document}
